% ─────────────────────────────────────────────────────────────────────────────
% SHARED coursework preamble — every course inputs THIS.
%
% PREAMBLE ONLY. No \begin{document}, no course-specific notation.
%
% Course-specific macros live at COURSE level, in
%   content/<course>/reference_docs/<course>_macros.tex
% which is inputted AFTER this file. ONE per course, shared by hw/qz/exam --
% macros under a single kind are unreachable from a sibling kind.
%
% That split is the point: the page style, the answer box, and the problem
% header are identical everywhere, so they are defined once; \conv and \ustep
% mean something in DSP and nothing in a networks course, so they are not.
%
% A per-problem file therefore starts:
%   % ─────────────────────────────────────────────────────────────────────────────
% SHARED coursework preamble — every course inputs THIS.
%
% PREAMBLE ONLY. No \begin{document}, no course-specific notation.
%
% Course-specific macros live at COURSE level, in
%   content/<course>/reference_docs/<course>_macros.tex
% which is inputted AFTER this file. ONE per course, shared by hw/qz/exam --
% macros under a single kind are unreachable from a sibling kind.
%
% That split is the point: the page style, the answer box, and the problem
% header are identical everywhere, so they are defined once; \conv and \ustep
% mean something in DSP and nothing in a networks course, so they are not.
%
% A per-problem file therefore starts:
%   % ─────────────────────────────────────────────────────────────────────────────
% SHARED coursework preamble — every course inputs THIS.
%
% PREAMBLE ONLY. No \begin{document}, no course-specific notation.
%
% Course-specific macros live at COURSE level, in
%   content/<course>/reference_docs/<course>_macros.tex
% which is inputted AFTER this file. ONE per course, shared by hw/qz/exam --
% macros under a single kind are unreachable from a sibling kind.
%
% That split is the point: the page style, the answer box, and the problem
% header are identical everywhere, so they are defined once; \conv and \ustep
% mean something in DSP and nothing in a networks course, so they are not.
%
% A per-problem file therefore starts:
%   \input{../../../../docs/latex/coursework_preamble.tex}
%   \input{../../reference_docs/cesc470_macros.tex}
%   \begin{document}
%
% Build-check one file, or a whole assignment:
%   docs/latex/build_tex.sh content/cesc_470/hw/hw01/p01_five_components.tex
%   docs/latex/build_tex.sh content/cesc_470/hw/hw01
%
% Doctrine: docs/directives/coursework-solutions.md
% ─────────────────────────────────────────────────────────────────────────────

\documentclass[11pt]{article}

\usepackage[margin=1in]{geometry}
\usepackage{amsmath, amssymb, mathtools}
\usepackage{graphicx}
\usepackage{float}
\usepackage{booktabs}
\usepackage{longtable}
\usepackage{enumitem}
\usepackage{siunitx}
\usepackage[dvipsnames]{xcolor}
\usepackage{tcolorbox}
\usepackage{fancyhdr}
\usepackage{caption}
\usepackage{tikz}
\usepackage{pgfplots}
\pgfplotsset{compat=1.18}
\usepackage[colorlinks=true, allcolors=NavyBlue]{hyperref}

\captionsetup{font=small, labelfont=bf}
\setlist[enumerate]{itemsep=2pt, topsep=4pt}
\setlist[itemize]{itemsep=2pt, topsep=4pt}

% --- final-answer box -------------------------------------------------------
% Every problem file ends with its result in one of these. The condensed
% solutions document is assembled by lifting exactly these.
\newtcolorbox{answerbox}[1][]{
  colback=Green!6, colframe=Green!55!black,
  boxrule=0.6pt, arc=2pt, left=6pt, right=6pt, top=4pt, bottom=4pt,
  #1
}

% --- citation to course material --------------------------------------------
% \cite{Module 01, PDF p55}  ->  a small italic parenthetical.
% Solutions cite the SLIDE/PAGE a formula came from so a grader (or a future
% reader) can check the claim against what was actually taught. See the
% directive: every non-obvious step carries one.
\newcommand{\srcref}[1]{{\small\itshape(#1)}}

% --- a step that came from OUTSIDE the course materials ---------------------
% Used only where the course is genuinely silent and the problem asks for
% outside research. Marked visually so it is never mistaken for lecture content.
\newcommand{\extref}[1]{{\small\itshape\color{Sepia}[external: #1]}}

% --- callout for the trap in a problem --------------------------------------
\newtcolorbox{trap}[1][]{
  colback=BurntOrange!7, colframe=BurntOrange!70!black,
  boxrule=0.6pt, arc=2pt, left=6pt, right=6pt, top=4pt, bottom=4pt,
  fonttitle=\bfseries, title=Watch out, #1
}

% --- per-problem header -----------------------------------------------------
% \problemheader{8}{15 pts}{Target clock rate}
% The optional 4-argument form carries a learning outcome where the course
% states one (CESC 410 does; CESC 470's HW1 does not).
\newcommand{\problemheader}[3]{%
  \begin{center}\large\bfseries
    Problem #1 \quad\normalsize\mdseries(#2)\\[2pt]
    \large\bfseries #3
  \end{center}
  \vspace{2pt}\hrule\vspace{8pt}
}
\newcommand{\problemheaderlo}[4]{%
  \begin{center}\large\bfseries
    Problem #1 \quad\normalsize\mdseries(#2, #3)\\[2pt]
    \large\bfseries #4
  \end{center}
  \vspace{2pt}\hrule\vspace{8pt}
}

% Course name in the running head. Defined BEFORE \lhead references it, and
% \renewcommand'd by each course's macros file (inputted after this one).
\providecommand{\coursename}{}

\pagestyle{fancy}
\fancyhf{}
\lhead{\small\coursename}
\rhead{\small\thepage}
\renewcommand{\headrulewidth}{0.3pt}

%   % CESC 470 Computer Architecture — course-specific macros ONLY.
%
% Inputted AFTER docs/latex/coursework_preamble.tex, which carries everything
% shared (answerbox, problemheader, srcref, page style). Put nothing general
% here: if another course would want it, it belongs in the shared preamble.

\renewcommand{\coursename}{CESC 470}

% --- performance-equation notation ------------------------------------------
% The course writes these in words on the slides; these render them compactly
% and, more importantly, IDENTICALLY across every problem file.
\newcommand{\CPUtime}{T_{\mathrm{CPU}}}
\newcommand{\IC}{\mathrm{IC}}                    % instruction count
\newcommand{\CPI}{\mathrm{CPI}}
\newcommand{\cycletime}{t_{\mathrm{cycle}}}
\newcommand{\clockrate}{f_{\mathrm{clk}}}
\newcommand{\cycles}{N_{\mathrm{cycles}}}
\newcommand{\speedup}{S}

% Amdahl's law, in the form the slides state it (Module 01, PDF p86, slide 49).
\newcommand{\amdahl}[3]{#1 + \dfrac{#2}{#3}}

%   \begin{document}
%
% Build-check one file, or a whole assignment:
%   docs/latex/build_tex.sh content/cesc_470/hw/hw01/p01_five_components.tex
%   docs/latex/build_tex.sh content/cesc_470/hw/hw01
%
% Doctrine: docs/directives/coursework-solutions.md
% ─────────────────────────────────────────────────────────────────────────────

\documentclass[11pt]{article}

\usepackage[margin=1in]{geometry}
\usepackage{amsmath, amssymb, mathtools}
\usepackage{graphicx}
\usepackage{float}
\usepackage{booktabs}
\usepackage{longtable}
\usepackage{enumitem}
\usepackage{siunitx}
\usepackage[dvipsnames]{xcolor}
\usepackage{tcolorbox}
\usepackage{fancyhdr}
\usepackage{caption}
\usepackage{tikz}
\usepackage{pgfplots}
\pgfplotsset{compat=1.18}
\usepackage[colorlinks=true, allcolors=NavyBlue]{hyperref}

\captionsetup{font=small, labelfont=bf}
\setlist[enumerate]{itemsep=2pt, topsep=4pt}
\setlist[itemize]{itemsep=2pt, topsep=4pt}

% --- final-answer box -------------------------------------------------------
% Every problem file ends with its result in one of these. The condensed
% solutions document is assembled by lifting exactly these.
\newtcolorbox{answerbox}[1][]{
  colback=Green!6, colframe=Green!55!black,
  boxrule=0.6pt, arc=2pt, left=6pt, right=6pt, top=4pt, bottom=4pt,
  #1
}

% --- citation to course material --------------------------------------------
% \cite{Module 01, PDF p55}  ->  a small italic parenthetical.
% Solutions cite the SLIDE/PAGE a formula came from so a grader (or a future
% reader) can check the claim against what was actually taught. See the
% directive: every non-obvious step carries one.
\newcommand{\srcref}[1]{{\small\itshape(#1)}}

% --- a step that came from OUTSIDE the course materials ---------------------
% Used only where the course is genuinely silent and the problem asks for
% outside research. Marked visually so it is never mistaken for lecture content.
\newcommand{\extref}[1]{{\small\itshape\color{Sepia}[external: #1]}}

% --- callout for the trap in a problem --------------------------------------
\newtcolorbox{trap}[1][]{
  colback=BurntOrange!7, colframe=BurntOrange!70!black,
  boxrule=0.6pt, arc=2pt, left=6pt, right=6pt, top=4pt, bottom=4pt,
  fonttitle=\bfseries, title=Watch out, #1
}

% --- per-problem header -----------------------------------------------------
% \problemheader{8}{15 pts}{Target clock rate}
% The optional 4-argument form carries a learning outcome where the course
% states one (CESC 410 does; CESC 470's HW1 does not).
\newcommand{\problemheader}[3]{%
  \begin{center}\large\bfseries
    Problem #1 \quad\normalsize\mdseries(#2)\\[2pt]
    \large\bfseries #3
  \end{center}
  \vspace{2pt}\hrule\vspace{8pt}
}
\newcommand{\problemheaderlo}[4]{%
  \begin{center}\large\bfseries
    Problem #1 \quad\normalsize\mdseries(#2, #3)\\[2pt]
    \large\bfseries #4
  \end{center}
  \vspace{2pt}\hrule\vspace{8pt}
}

% Course name in the running head. Defined BEFORE \lhead references it, and
% \renewcommand'd by each course's macros file (inputted after this one).
\providecommand{\coursename}{}

\pagestyle{fancy}
\fancyhf{}
\lhead{\small\coursename}
\rhead{\small\thepage}
\renewcommand{\headrulewidth}{0.3pt}

%   % CESC 470 Computer Architecture — course-specific macros ONLY.
%
% Inputted AFTER docs/latex/coursework_preamble.tex, which carries everything
% shared (answerbox, problemheader, srcref, page style). Put nothing general
% here: if another course would want it, it belongs in the shared preamble.

\renewcommand{\coursename}{CESC 470}

% --- performance-equation notation ------------------------------------------
% The course writes these in words on the slides; these render them compactly
% and, more importantly, IDENTICALLY across every problem file.
\newcommand{\CPUtime}{T_{\mathrm{CPU}}}
\newcommand{\IC}{\mathrm{IC}}                    % instruction count
\newcommand{\CPI}{\mathrm{CPI}}
\newcommand{\cycletime}{t_{\mathrm{cycle}}}
\newcommand{\clockrate}{f_{\mathrm{clk}}}
\newcommand{\cycles}{N_{\mathrm{cycles}}}
\newcommand{\speedup}{S}

% Amdahl's law, in the form the slides state it (Module 01, PDF p86, slide 49).
\newcommand{\amdahl}[3]{#1 + \dfrac{#2}{#3}}

%   \begin{document}
%
% Build-check one file, or a whole assignment:
%   docs/latex/build_tex.sh content/cesc_470/hw/hw01/p01_five_components.tex
%   docs/latex/build_tex.sh content/cesc_470/hw/hw01
%
% Doctrine: docs/directives/coursework-solutions.md
% ─────────────────────────────────────────────────────────────────────────────

\documentclass[11pt]{article}

\usepackage[margin=1in]{geometry}
\usepackage{amsmath, amssymb, mathtools}
\usepackage{graphicx}
\usepackage{float}
\usepackage{booktabs}
\usepackage{longtable}
\usepackage{enumitem}
\usepackage{siunitx}
\usepackage[dvipsnames]{xcolor}
\usepackage{tcolorbox}
\usepackage{fancyhdr}
\usepackage{caption}
\usepackage{tikz}
\usepackage{pgfplots}
\pgfplotsset{compat=1.18}
\usepackage[colorlinks=true, allcolors=NavyBlue]{hyperref}

\captionsetup{font=small, labelfont=bf}
\setlist[enumerate]{itemsep=2pt, topsep=4pt}
\setlist[itemize]{itemsep=2pt, topsep=4pt}

% --- final-answer box -------------------------------------------------------
% Every problem file ends with its result in one of these. The condensed
% solutions document is assembled by lifting exactly these.
\newtcolorbox{answerbox}[1][]{
  colback=Green!6, colframe=Green!55!black,
  boxrule=0.6pt, arc=2pt, left=6pt, right=6pt, top=4pt, bottom=4pt,
  #1
}

% --- citation to course material --------------------------------------------
% \cite{Module 01, PDF p55}  ->  a small italic parenthetical.
% Solutions cite the SLIDE/PAGE a formula came from so a grader (or a future
% reader) can check the claim against what was actually taught. See the
% directive: every non-obvious step carries one.
\newcommand{\srcref}[1]{{\small\itshape(#1)}}

% --- a step that came from OUTSIDE the course materials ---------------------
% Used only where the course is genuinely silent and the problem asks for
% outside research. Marked visually so it is never mistaken for lecture content.
\newcommand{\extref}[1]{{\small\itshape\color{Sepia}[external: #1]}}

% --- callout for the trap in a problem --------------------------------------
\newtcolorbox{trap}[1][]{
  colback=BurntOrange!7, colframe=BurntOrange!70!black,
  boxrule=0.6pt, arc=2pt, left=6pt, right=6pt, top=4pt, bottom=4pt,
  fonttitle=\bfseries, title=Watch out, #1
}

% --- per-problem header -----------------------------------------------------
% \problemheader{8}{15 pts}{Target clock rate}
% The optional 4-argument form carries a learning outcome where the course
% states one (CESC 410 does; CESC 470's HW1 does not).
\newcommand{\problemheader}[3]{%
  \begin{center}\large\bfseries
    Problem #1 \quad\normalsize\mdseries(#2)\\[2pt]
    \large\bfseries #3
  \end{center}
  \vspace{2pt}\hrule\vspace{8pt}
}
\newcommand{\problemheaderlo}[4]{%
  \begin{center}\large\bfseries
    Problem #1 \quad\normalsize\mdseries(#2, #3)\\[2pt]
    \large\bfseries #4
  \end{center}
  \vspace{2pt}\hrule\vspace{8pt}
}

% Course name in the running head. Defined BEFORE \lhead references it, and
% \renewcommand'd by each course's macros file (inputted after this one).
\providecommand{\coursename}{}

\pagestyle{fancy}
\fancyhf{}
\lhead{\small\coursename}
\rhead{\small\thepage}
\renewcommand{\headrulewidth}{0.3pt}

% CESC 470 Computer Architecture — course-specific macros ONLY.
%
% Inputted AFTER docs/latex/coursework_preamble.tex, which carries everything
% shared (answerbox, problemheader, srcref, page style). Put nothing general
% here: if another course would want it, it belongs in the shared preamble.

\renewcommand{\coursename}{CESC 470}

% --- performance-equation notation ------------------------------------------
% The course writes these in words on the slides; these render them compactly
% and, more importantly, IDENTICALLY across every problem file.
\newcommand{\CPUtime}{T_{\mathrm{CPU}}}
\newcommand{\IC}{\mathrm{IC}}                    % instruction count
\newcommand{\CPI}{\mathrm{CPI}}
\newcommand{\cycletime}{t_{\mathrm{cycle}}}
\newcommand{\clockrate}{f_{\mathrm{clk}}}
\newcommand{\cycles}{N_{\mathrm{cycles}}}
\newcommand{\speedup}{S}

% Amdahl's law, in the form the slides state it (Module 01, PDF p86, slide 49).
\newcommand{\amdahl}[3]{#1 + \dfrac{#2}{#3}}


% CONDENSED ANSWER SHEET. Final answers only -- no derivations.
% Assembled by lifting the answerbox from each per-problem file:
%   p01_five_components  p02_von_neumann_tradeoff  p03_instruction_encodings
%   p04_architecture_vs_organization  p05_same_isa_different_org
%   p06_isa_components  p07_arm_vs_x86  p08_target_clock_rate
%   p09_cpi_and_execution_time  p10_amdahl_speedup
% Doctrine: docs/directives/coursework-solutions.md

\title{\vspace{-1.5cm}CESC 470 --- Homework 1 Solutions\\
       {\large Computer Architecture}}
\author{}
\date{Fall 2026 \quad\textbullet\quad Due 9/13/2026 \quad\textbullet\quad 100 points}

\begin{document}
\maketitle
\vspace{-1.0cm}

%---------------------------------------------------------------------------
\section*{1. Five classic components \normalsize(5 pts)}

\begin{enumerate}[nosep, leftmargin=1.6em]
  \item \textbf{Control unit} — fetches and decodes instructions, issues the
    control signals that execute them.
  \item \textbf{Datapath} — the ALU and registers; performs the computation.
  \item \textbf{Memory} — stores both program and data.
  \item \textbf{Input} — brings data into the machine.
  \item \textbf{Output} — sends results out.
\end{enumerate}
Control unit $+$ datapath $=$ the \textbf{CPU}; all five joined by a common
address/data/control \textbf{bus}.

%---------------------------------------------------------------------------
\section*{2. Von Neumann trade-off \normalsize(10 pts)}

\textbf{Advantage — general-purpose reprogrammability.} Instructions are bit
sequences in ordinary read/write memory, so programs are \emph{loaded} rather
than wired in, and programs (compilers, assemblers, loaders) can write other
programs.

\textbf{Disadvantage — the von Neumann bottleneck.} Instructions and data share
one memory and one bus, so instruction fetch competes with data access every
cycle and memory bandwidth caps performance regardless of CPU speed. Writable
instruction memory also lets code be overwritten or data executed. The
\emph{Harvard} architecture avoids this with separate memories, losing the
flexibility above.

%---------------------------------------------------------------------------
\section*{3. Unique 32-bit instructions \normalsize(5 pts)}

$$2^{32} = 4{,}294{,}967{,}296 \approx 4.29\times10^{9}$$
Each of the 32 bit positions independently takes 2 values. This is an
\emph{upper bound} on bit patterns; a real ISA splits the 32 bits into opcode
and operand fields, so far fewer distinct \emph{operations} are encoded.

%---------------------------------------------------------------------------
\section*{4. Architecture vs organization \normalsize(10 pts)}

\textbf{Architecture tells you \emph{what} a computer does; organization tells
you \emph{how} it does it.}

\textbf{Architecture} $=$ \textbf{ISA $+$ hardware organization}. The ISA is the
abstract, programmer-visible programming model — instruction set, register set,
addressing modes — and the assembly language embodies it.

\textbf{Organization} (today: \emph{microarchitecture}) is the implementation of
an ISA: pipeline depth, caches, clock rate, buses. Any ISA admits many
organizations, and vendors change the organization every generation while
holding the ISA constant.

\emph{Test:} if a program can \textbf{detect} it, architecture; if it only
changes how \textbf{fast} the same program runs, organization.

%---------------------------------------------------------------------------
\section*{5. Same ISA, different organizations \normalsize(5 pts)}

A \textbf{desktop and a laptop x86-64 processor}. Both implement the same ISA, so
both run the \emph{same compiled binary} with identical results. Their
organizations differ — clock rate, core count, cache sizes, pipeline width,
power budget — and the program can detect the difference only as \emph{run
time}.

%---------------------------------------------------------------------------
\section*{6. Three components of an ISA \normalsize(3 pts)}

\begin{enumerate}[nosep, leftmargin=1.6em]
  \item \textbf{Instruction set (opcodes)} — the operations and their encodings.
  \item \textbf{Organization of storage / register set} — number, width, purpose
    of programmer-visible registers.
  \item \textbf{Addressing modes} — how an instruction names its operand.
\end{enumerate}
{\small(Also valid: data types; instruction encoding; program-visible exception
handling. All are \emph{programmer-visible} — caches and clock rate are
organization, not ISA.)}

%---------------------------------------------------------------------------
\section*{7. ARM vs x86 \normalsize(25 pts)}

\begin{center}
% Middle column at 3.0cm hyphenated "Power con-sumption" / "Decode com-plexity";
% 3.9cm fits both on one line. Last column reduced to keep the total width.
\begin{tabular}{@{}p{0.5cm}p{3.9cm}p{10.3cm}@{}}
  \toprule
  \multicolumn{3}{@{}l}{\textbf{ARM}} \\
  \midrule
  $+$ & \textbf{Power consumption}
      & Fixed-width, regular RISC encoding keeps the decoder small and simple,
        so less energy per instruction. Hence its dominance in battery-powered
        devices, where performance-per-watt is the metric. \\
  $-$ & \textbf{Compatibility}
      & The large body of existing desktop/server software was compiled for
        x86 and must be recompiled or emulated (Rosetta 2, Windows-on-ARM),
        costing performance and risking edge cases. \\
  \midrule
  \multicolumn{3}{@{}l}{\textbf{x86}} \\
  \midrule
  $+$ & \textbf{Compatibility}
      & Binary compatibility maintained since the 8086 (1978); decades-old
        software still runs on current processors — decisive for legacy and
        closed-source codebases. \\
  $-$ & \textbf{Decode complexity / power}
      & Variable-length instructions (1--15 bytes) mean instruction $n+1$'s
        start is unknown until $n$ is decoded. Parallel decode needs
        substantial extra logic, and cores translate to internal micro-ops —
        transistors and power spent servicing a legacy encoding that cannot be
        dropped without losing the advantage above. \\
  \bottomrule
\end{tabular}
\end{center}

{\small\textbf{Caveat:} raw performance is governed mainly by \emph{organization}
(pipeline, caches, process node), not the ISA. The durable ISA-level differences
are decode complexity and energy per instruction.}

%---------------------------------------------------------------------------
\section*{8. Target clock rate for Computer B \normalsize(15 pts)}

\begin{align*}
  \cycles^{A} &= 18\ \text{s} \times 600\ \text{MHz} = 1.08\times10^{10} \\
  \cycles^{B} &= 1.5 \times 1.08\times10^{10} = 1.62\times10^{10}
     \qquad \CPUtime^{B} = \tfrac{18}{2} = 9\ \text{s} \\
  \clockrate^{B} &= \frac{1.62\times10^{10}}{9\ \text{s}}
     = \mathbf{1.8\ \text{GHz}}
\end{align*}
Equivalently $600\ \text{MHz} \times 1.5 \times 2$ — the clock must cover
\emph{both} the $2\times$ speed target and the $1.5\times$ cycle penalty.
\emph{(1.2 GHz is the trap: it ignores the extra cycles.)}

%---------------------------------------------------------------------------
\section*{9. Average CPI and execution time \normalsize(12 pts)}

\textbf{(a)} $\;\CPI = (0.40)(1) + (0.40)(2) + (0.20)(3) = 0.4+0.8+0.6
  = \mathbf{1.8}$ cycles/instruction.

\textbf{(b)} $\;\CPUtime = \IC \times \CPI \times \cycletime
  = (10^{6})(1.8)(1\ \text{ns}) = 1.8\times10^{6}\ \text{ns}
  = \mathbf{1.8\ \text{ms}}$.

%---------------------------------------------------------------------------
\section*{10. Amdahl's law \normalsize(10 pts)}

$$\CPUtime^{\text{new}} = 0.20 + \frac{0.80}{10} = 0.28
\qquad\Longrightarrow\qquad
\speedup = \frac{1}{0.28} = \mathbf{3.57\times}$$

\textbf{Theoretical ceiling} (infinitely fast FP): $\;\speedup_{\max}
= 1/0.20 = \mathbf{5\times}$.

{\small The stated $10\times$ improvement achieves 3.57 of the 5 possible in the
limit. \emph{``Maximum possible'' is ambiguous} — 3.57$\times$ answers the
improvement given; 5$\times$ is the ceiling. Both are stated.}

\end{document}
